\chapter{The test and its economics}
\label{ch:evaluation}

The product changes a chain of work: an author revises, a reader reads, and a
sponsor decides whether the result is worth keeping. Document scores observe
only one part of that chain, so the evaluation compares the complete workflow.
The central question is:

\begin{quote}
Does the Humanvoice workflow help more documents reach a reader decision within
a fixed horizon, or reduce the attention required to get there, without
increasing meaning-preservation failures relative to the incumbent bundle?
\end{quote}

One live document cannot answer the causal question. It can show that the
system runs, that people can use its records, and that the team can observe
burden and failures. An efficacy estimate requires a later comparison across
documents and workflows.

Figure~\ref{fig:evaluation-logic} marks the decision boundary. The live case
can expose integration failures and burden; only a later comparison can
estimate an effect.

\hvFigureEvaluationLogic

\section{The estimand}

Let $i$ index a document task, and let $w=0$ denote the incumbent bundle of
ordinary linting, general LLM editing, version control, citation checking, and
manual review. Let $w=1$ denote the Humanvoice workflow. Before assignment, the
study fixes a maximum of $K$ review cycles and $H$ calendar days. At that
horizon it records one disposition: accepted, revision still requested,
declined, withdrawn by the author, or incomplete for an operational reason.

Define $A_i^H(w)=1$ when the named reader accepts by the horizon,
$T_i^H(w)$ as cumulative reader hours through the horizon,
$R_i^K(w)=\min(R_i(w),K)$ as restricted review cycles, and $M_i(w)$ as an
indicator that a meaning-bearing object changed without an explained author
decision. The co-primary estimands are
\begin{equation}
  \Delta_A=\mathbb E[A_i^H(1)-A_i^H(0)],
  \qquad
  \Delta_T=\mathbb E[T_i^H(1)-T_i^H(0)].
  \label{eq:estimand}
\end{equation}

Acceptance prevents a quick decline from looking like success; restricted
reader time captures the scarce attention the product is meant to save.
$R_i^K$ remains an important process outcome, reported with the terminal
disposition rather than only for accepted documents. Every decline,
withdrawal, and incomplete task stays in the analysis. An unexplained protected
change is a safety failure regardless of time or acceptance. Author editing
time, false-positive burden, substantive revision rate, reader reconstruction,
confidence, and time to decision explain how the primary outcomes arose.

\section{The unit and the comparator}

Give each task a named reader and a decision before assigning a workflow. The
task may be a funding proposal, model note, technical report, or manuscript, but
its genre and decision must be recorded. Both conditions receive the same
source snapshot and protected-object inventory. The reader should be blinded to
the workflow where practical; the author cannot be, so the study measures
author expectations rather than pretending they do not exist.

The comparator is the practice the team would use if Humanvoice did not exist,
not an ideal editor invented for the experiment. Before the study begins, write
down which linter and general editor people use, how they check version changes
and citations, who reads, and what counts as acceptance. A product that wins
against an artificially weak baseline has not earned its place.

\section{The reader record}

The reader rubric is short enough to use repeatedly and specific enough to
distinguish comprehension from agreement:

\begin{table}[H]
\centering
\small
\begin{tabularx}{\textwidth}{@{}p{0.28\textwidth}L p{0.25\textwidth}@{}}
\toprule
Question & Response recorded & Why it matters \\
\midrule
What problem or decision is this document about? & Free response plus
confidence. & Tests whether purpose is visible. \\
What is the proposed mechanism or object? & Free response plus selected
components. & Tests whether the argument has an identifiable center. \\
Which evidence is decisive? & Source or passage selection with explanation. &
Tests evidence-to-conclusion linkage. \\
What qualification changes your interpretation? & Short response. & Tests
whether uncertainty is visible without swallowing the argument. \\
What action would you take now? & Accept, request revision, or decline, with
reason. & Connects comprehension to the actual decision. \\
Where did you stop or reread? & Page or section, a short quoted span, and the
missing relation you were trying to recover. & Locates the reader's first
failure instead of reducing it to a score. \\
\bottomrule
\end{tabularx}
\caption{A reader outcome records reconstruction and action, not stylistic
approval.}
\label{tab:reader-rubric}
\end{table}

Reader disagreement is not automatically noise. If two qualified readers
recover different mechanisms, that may indicate an ambiguous document. The
analysis should report agreement, confidence, and the reasons for requested
changes rather than force a single hidden "quality" label.

\section{The feasibility case}

The first twelve weeks include one live document. Its purpose is integration,
not efficacy. The case must demonstrate:

\begin{enumerate}
\item the source can be parsed and rendered with protected spans intact;
\item findings are reproducible and actionable in the author's editor;
\item accepted and rejected revisions produce an understandable protected
      comparison;
\item privacy and disclosure choices are respected; and
\item a reader can complete the rubric without repository context.
\end{enumerate}

The case should be selected because its failure cost is real and its owner is
willing to permit a reading by someone unfamiliar with the project, not because
it is likely to make the system look good. The result note reports all abstentions, unresolved
citations, reader requests, and meaning checks. It cannot report a reduction in
rounds relative to a counterfactual that was never run.

\section{The comparative study}

After the feasibility case, the evaluation lead should freeze a corpus with
tuning and held-out partitions. The corpus needs documents from more than one
writer and more than one task class. If the available volume is small, the
study should say so and use paired designs or repeated reader tasks rather than
inventing a large-sample certainty.

The study should report outcomes by baseline writer experience, document genre,
reader role, and use of model assistance. Averages can conceal harm to the
authors or readers whose work matters most. The field literature's
heterogeneous effects make this stratification a design requirement rather than
an optional subgroup analysis \citep{brynjolfsson2025work,dellacqua2026jagged}.

Model-based screening, if used, is calibrated on a frozen human-labelled set.
Pair order is randomized, length is controlled, and the judge family is varied
when feasible. A model result can nominate a passage or explain a discrepancy;
it cannot accept a document.

\section{The economic decision rule}

One accounting identity is enough. Let $N$ be annual in-scope document volume,
$T_w$ cumulative reader hours per document through the fixed horizon, $W_w$
author and operating hours, $c_r$ and $c_a$ their loaded rates, $p_wL_w$ the
expected cost of a meaning incident, and $C$ annualized product cost. The
annual net value of replacing workflow 0 with workflow 1 is
\begin{equation}
  V=N\left[(T_0-T_1)c_r+(W_0-W_1)c_a
    -(p_1L_1-p_0L_0)\right]-C.
  \label{eq:value}
\end{equation}

This is an accounting model, not a forecast. Acceptance by the horizon is a
condition on using it: a workflow that saves time by producing more declines or
unfinished documents does not receive a favorable value estimate. The
feasibility run adds the quantities needed to populate the model:

\begin{description}
\item[Volume.] How many in-scope documents does the group produce each year?
\item[Baseline burden.] How many rounds and reader hours do those documents
  consume, with a distribution rather than one anecdotal average?
\item[Workflow burden.] How much author and reviewer time does Humanvoice add,
  including false-positive resolution and protected comparisons?
\item[Loaded rates.] What is the local cost of author, reader, evaluator, and
  engineer time?
\item[Build and maintenance.] What person-weeks and recurring infrastructure
  cost are required after the pilot?
\end{description}

Until those inputs exist, symbolic break-even cases are more honest than
invented volumes and hourly rates.

\begin{table}[H]
\centering
\small
\begin{tabularx}{\textwidth}{@{}p{0.26\textwidth}p{0.34\textwidth}L@{}}
\toprule
Decision question & Break-even expression & Observation required \\
\midrule
Maximum annual product cost &
$C^*=N[(T_0-T_1)c_r+(W_0-W_1)c_a-(p_1L_1-p_0L_0)]$ &
Local volume, time distributions, loaded rates, and incident assumptions. \\
Reader-time saving required when all other changes are zero &
$T_0-T_1=C/(Nc_r)$ & Canonical reader-session timing under both workflows. \\
Maximum extra author time the workflow can consume &
$W_1-W_0=[(T_0-T_1)c_r-(p_1L_1-p_0L_0)-C/N]/c_a$ &
Author event logs and adjudicated false-alert work. \\
Effect of one additional meaning risk &
$\Delta V=-N(p_1L_1-p_0L_0)$ & Incident definition, probability range, and
loss range approved by the sponsor. \\
\bottomrule
\end{tabularx}
\caption{Symbolic sensitivity cases. Each row shows which local observation,
not a placeholder forecast, determines the investment boundary.}
\label{tab:value-sensitivity}
\end{table}

The sponsor should see low, central, and high cases for every uncertain input,
with each value marked observed, assumed, or open. A point estimate made before
the baseline exists would be financial decoration rather than analysis.

\section{Stop and redirect conditions}

The evaluation has four stop conditions:

\begin{enumerate}
\item the protected comparison misses or silently changes a meaning-bearing
      object;
\item the diagnostic queue consumes more author or reader time than it saves;
\item the privacy or disclosure boundary cannot be implemented for the live
      document; or
\item the reader cannot recover the purpose and requested action after two
      bounded repair cycles.
\end{enumerate}

A failed candidate does not automatically reject the product direction. If a
parser fails malformed macros but a documented fallback can repair that exact
failure, the next phase tests the fallback. If the reader endpoint itself is
not observable or the burden cannot be measured, the direction is not ready for
a larger study. The result note must distinguish those cases.
This feasibility case establishes only that the workflow can be run and measured
on one live document. It cannot establish a general reduction in revision
rounds, a universal improvement in technical prose, or a return on investment;
those claims require the comparative study and its uncertainty.
